Тем не менее, секвенциальная топология сохраняет некоторые свойства изначальной топологии.
В частности, верно следующее утверждение.
\begin{proposition}
    Пусть $(X, \tau) \in \cat{Top}$ и $\tau_{seq}$ -- секвенциальная топология.
    Тогда для всяких $\{x_n\} \subset X$ и $x \in X$ верно следующее:
    \[
        x_n \xrightarrow{\tau} x \Longleftrightarrow x_n \xrightarrow{\tau_{seq}} x.
    \]
\end{proposition}
\begin{proof}
    Очевидно, что если мы сходимся по $\tau_{seq}$, то мы сходимся и по $\tau$ -- поскольку $\tau \subset \tau_{seq}$, и тогда, по определению:
    \[
        \forall U(x) \in \tau_{seq} \ \exists N \in \N\colon \forall n \geq N \hookrightarrow x_n \in U(x),
    \]
    а следовательно, так как $\tau \subset \tau_{seq}$, то и
    \[
        \forall U(x) \in \tau \ \exists N \in \N \colon \forall n \geq N \hookrightarrow x_n \in U(x),
    \]
    и $x_n \xrightarrow{\tau} x$.

    Покажем обратную импликацию.
    Пусть $x_n \xrightarrow{\tau} x$ и, предположим, что $x_n \xcancelrightarrow{\tau_{seq}} x$.
    Отсутствие сходимости по $\tau_{seq}$ обозначает что $\exists U(x) \in \tau_{seq}$ такое, что
    \[
        \forall N \in \N \ \exists n \in \N\colon n \geq N \hookrightarrow x_n \notin U(x).
    \]
    Тогда у нас существует подпоследовательность $\{x_{n_k}\} \subset X \setminus U(x)$.
    Но она всё также сходится к $x\colon x_{n_k} \xrightarrow{\tau} x$.
    Поскольку $U(x) \in \tau_{seq}$, то его дополнение секвенциально замкнуто и, в частности, предел $x_{n_k}$ лежит в $X \setminus U(x)$ -- противоречие.
\end{proof}
\begin{note}
    Это утверждение, в частности, показывает, что $S \subset X$ замкнуто в $(X, \tau_{seq})$ тогда и только тогда, когда оно секвенциально замкнуто в $(X, \tau_{seq})$.

    Распишем это подробнее.
    
    Достаточно показать, что $\forall S \subset X$ -- секвенциально-замкнутое в $(X, \tau_{seq})$ будет являться и замкнутым в $(X, \tau_{seq})$.
    Но это эквивалентно тому, что $S \in \phi_{seq}$, то есть тому, что $S$ -- секвенциально-замкнуто в $(X, \tau)$.

    А это очевидно в силу эквивалентности сходимостей.
\end{note}
\begin{note}
    Заметим, что при этом в $(X, \tau_{seq})$ всё ещё может оставаться отличие между замыканием и секвенциальным замыканием.
    Замыкание в $(X, \tau_{seq})$ имеет вид:
    \[
        [S]_{\tau_{seq}} = \bigcap_{S \subset A \in \phi_{seq}} A.
    \]
    Секвенциальное замыкание остаётся тем же.
\end{note}

\section{Аксиома счётности (первая).}
\begin{definition}
    Пусть $(X, \tau) \in \cat{Top}$ и $x \in X$.
    Семейство $\beta(x) \subset \tau\colon x \in V(x) \in \beta(x) \subset \tau$ будем называть локальной базой (топологии $\tau$ в точке $x$), если $\forall U(x) \in \tau$ -- окрестности точки $x$ существует $V \in \beta(x)$ такое, что $V \subset U(x)$.
\end{definition}
\begin{definition}
    Пусть $(X, \tau) \in \cat{Top}$.
    Будем говорить что оно удовлетворяет аксиоме счётности (на самом деле -- первой аксиоме счётности\footnote{Топологическое пространство удовлетворяет второй аксиоме счётности, если оно имеет счётную базу топологии -- определение базы топологии будет дальше.})
    если у каждой точки $x \in X$ существует не более чем счётная локальная база окрестностей $\beta(x) = \{V_n\}$.
\end{definition}
\begin{note}
    Поскольку мы можем рассмотреть семейство $\{W_n\}$, где
    \[
        W_n = \bigcap_{k = 1}^{n} V_k,
    \]
    и $\{W_n\}$ -- тоже будет не более чем счётной локальной базой окрестностей, то мы можем без ограничения общности считать локальную базу убывающей последовательностью множеств.
\end{note}
\begin{proposition}
    Пусть $(X, \tau) \in \cat{Top}$ удовлетворяет первой аксиоме счётности и $S \subset X$.
    Тогда любая топологическая точка прикосновения $S$ является и секвенциальной точкой прикосновения $S$.

    Также верно то, что $S$ -- секвенциально замкнуто $\Longleftrightarrow$ $S$ -- замкнуто и что $[S]_{seq} = [S]_{\tau}$.
\end{proposition}
\begin{proof}
    Второе утверждение тривиально следует из первого: поскольку всякое замкнутое является и секвенциально замкнутым, и, в случае справедливости первого утверждения, секвенциально замкнутое множество $S$ будет содержать все свои секвенциальные точки прикосновения -- следовательно оно будет содержать все топологические точки прикосновения и будет замкнутым.
    Равенство замыканий следует из эквивалентности понятий <<секвенциальная точка прикосновения>> и <<топологическая точка прикосновения>>.

    Покажем справедливость первого утверждения.
    Пусть $S \subset X$ и $x$ -- топологическая точка прикосновения $S$, то есть
    \[
        \forall U(x) \in \tau \hookrightarrow U(x) \cap S \neq \emptyset.
    \]
    В окрестности точки $x$ существует невозрастающая, не более чем счётная, локальная база $\beta(x) = \{W_n\}_{n = 1}^{\infty}$.
    Поскольку $\forall n \in \N \hookrightarrow W_n \in \tau$, то $\forall W_n \hookrightarrow W_n \cap S \neq \emptyset$.

    Пусть $x_n \in W_n \cap S$ -- такие точки существуют.
    Покажем, что $x_n \xrightarrow{\tau} x$.
    Пусть $U(x) \in \tau$.
    По определению локальной базы, существует $N \in \N\colon W_N \subset U(x)$.
    В силу упорядоченности, это верно и $\forall n \geq N\colon W_n \subset W_N \subset U(x)$.
    Следовательно $\forall n \geq N \hookrightarrow x_n \in W_n \subset U(x)$ -- что по определению и даёт сходимость последовательности к $x$.
    Поскольку $\{x_n\} \subset S$, то $x$ -- секвенциальная точка прикосновения $S$ по определению и утверждение доказано.
\end{proof}
\section{Сходимость, порождающая топологию.}
\begin{definition}
    Пусть $X$ -- некоторое непустое множество.
    Пусть на нём задана сходимость <<$\rightarrow$>> (то есть всякая последовательность $\{x_n\} \subset X$ объявляется либо сходящейся к некоторой $x \in X$, либо расходящейся\footnote{Кажется, формально это можно определить как отображение из $X^{\N}$ -- пространства последовательностей -- в $2^X$, где под $2^X$ понимается множество всевозможных пределов последовательности}), удовлетворяющая следующим свойствам:
    \begin{enumerate}
        \item Если $\{x_n\} \subset X$ и $x_n \rightarrow x$, то для всякой $\{x_{n_k}\}$ -- подпоследовательности $x_n$ -- верно, что $x_{n_k} \rightarrow x$.
        \item Если $\{x_n\}$ -- стационарная последовательность и стабилизируется на $x$, то $x_n \rightarrow x$.
    \end{enumerate}
\end{definition}
\begin{note}
    Для всякого пространства со сходимостью -- то есть пары $(X, \rightarrow)$ мы можем определить аналогичным образом секвенциальные точки прикосновения (то есть $x \in X$ является секвенциальной точкой прикосновения множества $S$, если $\exists \{x_n\}\subset S\colon x_n \rightarrow x$).
    Аналогичным образом мы можем определить секвенциально-замкнутые множества, как множества, содержащие все свои секвенциальные точки прикосновения.
    Пусть $\phi_{seq}$ -- семейство секвенциально-замкнутых множеств относительно сходимости $\rightarrow$.

    Понятно, что $\emptyset$ и $X$ лежат в $\phi_{seq}$.

    Покажем, что $\phi_{seq}$ замкнуто относительно произвольных пересечений.
    Пусть $\{W_\alpha\} \subset \phi_{seq}$ и $W = \bigcap_{\alpha} W_\alpha$.
    Если $x$ -- секвенциальная точка прикосновения $W$, то, в частности, она является и секвенциальной точкой прикосновения всякого $W_\alpha \supset W$.
    Следовательно, $x \in W_\alpha$ при всяком $\alpha$, а значит лежит и в $W$.

    Покажем, что $\phi_{seq}$ замкнуто относительно конечных объединений.
    Пусть $\{W_k\}_{k = 1}^{N} \subset \phi_{seq}$ и $W = \bigcup_{k = 1}^{N} W_k$.
    Если $x$ -- секвенциальная точка прикосновения $W$, то существует последовательность $\{x_n\} \subset W$ такая, что $x_n \rightarrow x$.
    Поскольку $W = \bigcup_{k = 1}^{N} W_k$ то существует $j \in \{1, \ldots, N\}$ и некоторая подпоследовательность $\{x_{n_k}\} \subset W_j$.
    По определению сходимости $x_{n_j} \rightarrow x$.
    Следовательно, в силу секвенциальной замкнутости $W_j$, $x \in W_j$.
    Но тогда $x \in W$ -- и $W$ -- секвенциально замкнуто.

    Видно, что $\phi_{seq}$ -- семейство замкнутых множеств и порождает некоторую топологию $\tau_{seq}$ на $X$.
\end{note}
\begin{proposition}
    Пусть задано пространство со сходимостью $(X, \rightarrow)$.
    Тогда $\forall \{x_n\} \subset X$ верна следующая импликация:
    \[
        x_n \rightarrow x \Rightarrow x_n \xrightarrow{\tau_{seq}} x.
    \]
    Обратное, в общем случае, неверно.
\end{proposition}
\begin{proof}
    Пусть $x_n \rightarrow x$, но не сходится к $x$ по $\tau_{seq}$.
    Тогда $\exists U(x) \in \tau_{seq}$ такая, что
    \[
        \forall N \in \N \exists n > N\colon x_n \notin U(x).
    \]
    Следовательно, существует подпоследовательность $\{x_{n_k}\} \subset X \setminus U(x)$.
    Заметим, что по определению сходимости $x_{n_k} \rightarrow x$.
    Поскольку $X \setminus U(x)$ -- секвенциально замкнуто, то $x \in X \setminus U(x)$ -- противоречие.
\end{proof}
\section{Метрическая топология.}
\begin{example}
    Пусть $X$ -- множество.
    Функция $\rho\colon X \times X \to \R$
    называется метрикой, если выполняются следующие условия:
    \begin{enumerate}
        \item $\forall x, y \in X \hookrightarrow \rho(x, y) \geq 0$.
        \item $\forall x, y \in X \hookrightarrow \rho(x, y) = \rho(y, x)$.
        \item $\forall x, y, z \in X \hookrightarrow \rho(x, y) \leq \rho(x, z) + \rho(z, y)$.
        \item $\forall x, y \in X \hookrightarrow \rho(x, y) = 0 \Longleftrightarrow x = y$.
    \end{enumerate}
    Пара $(X, \rho)$ называется метрическим пространством.

    Сходимость в $(X, \rho)$ определяется как $\{x_n\} \subset X\colon$
    \[
        x_n \xrightarrow{\rho} x \Longleftrightarrow \rho(x_n, x) \rightarrow 0.
    \]

    Отсюда очевидным образом следует, что последовательность имеет не более чем один предел1, поскольку если $x_n \xrightarrow{\rho} x$, и $x_n \xrightarrow{\rho} y$, то
    \[
        \rho(x, y) \leq \rho(x, x_n) + \rho(y, x_n) \rightarrow 0 \Rightarrow \rho(x, y) = 0 \Rightarrow x = y.
    \]

    Очевидно, что стационарные последовательности сходятся и что всякая подпоследовательность сходящейся последовательности сходится к тому же элементу.
    Следовательно, $(X, \xrightarrow{\rho})$ является пространством со сходимостью и, в частности, порождает некоторую топологию $\tau_{seq}$.

    Ранее было показано, что при применении такой конструкции мы получаем следующую импликацию:
    \[
        x_n \xrightarrow{\rho} x \Rightarrow x_n \xrightarrow{\tau_{seq}} x.
    \]
    Покажем, что в этом случае верно её обращение.

    Пусть $r > 0$ и $x \in X$.
    Рассмотрим множество
    \[
        M(x, r) = \{y \in X \mid \rho(x, y) \geq r\}.
    \]
    Очевидно, что такое множество является секвенциально замкнутым, поскольку если $\{x_n\} \subset M(x, r)\colon x_n \xrightarrow{\rho} y \in X$, то
    \[
        \rho(x, y) \geq \rho(x, x_n) - \rho(x_n, y) \geq r - \rho(x_n, y) \rightarrow r.
    \]
    Значит $O(x, r) = X \setminus M(x, r)$ является открытым по этой топологии.
    Заметим, что $O(x, r)$ является окрестностью $x$.
    Тогда, если $x_n \xrightarrow{\tau_{seq}} x$, то
    \[
        \forall r > 0 \exists N \in \N  \forall n \geq N \hookrightarrow x_n \in O(x, r).
    \]
    Поэтому $\forall r > 0 \exists N \in \N \forall n \geq N \hookrightarrow \rho(x_n, x) \leq r$ и $x_n \xrightarrow{\rho} x$, а значит
    \[
        x_n \xrightarrow{\rho} x \Longleftrightarrow x_n \xrightarrow{\tau_{seq}} x.
    \]
    Пусть $\tau_\rho \stackrel{def}{=} \tau_{seq}$.
\end{example}